\section{Contributions and Provenance}

This project is largely assembly of existing work. The table below states what is
established, what is being replicated, and what would be new, so that no part of this
report implies more novelty than it has.

\begin{table}[h]
\centering
\small
\begin{tabular}{p{5.2cm} p{2.6cm} p{6.2cm}}
\toprule
\textbf{Component} & \textbf{Status} & \textbf{Source} \\
\midrule
JEPA architecture & Established & \citet{assran2023ijepa} \\
SIGReg loss & Established, not ours & \citet{lejepa2025}; verified line-by-line against the reference implementation, 2026-08-14 \\
EMA target + stop-gradient & Established & \citet{grill2020byol,chen2021simsiam} \\
Collapse diagnostics & Established practice & Standard in the SSL literature \\
Linear-probe protocol & Established practice & Standard SSL evaluation \\
Four-condition ablation grid & Standard methodology & Conventional design \\
Scale-invariance of standardized probing and cosine retrieval & Established by definition & Property of the metrics; predicted a priori, not discovered here \\
RankMe & Established, not ours & \citet{garrido2023rankme} \\
Scale-dependence of stop-gradient under SIGReg & \textbf{Replication attempt} & Encountered secondhand; primary source not read end-to-end \\
Cheap diagnostics as early warning & \textbf{Tested, not supported} & Phase-1 pilot: the probe turned first (Section~\ref{sec:results}) \\
\emph{Magnitude} of the blindness: a collapsed encoder scoring 0.413 on the standardized probe and $1.8\times$ chance on retrieval & \textbf{Ours, claim-grade} & Phase~A, five paired seeds, paired Student-t (Section~\ref{sec:results}); probe endpoints provisional, see Section~\ref{sec:provenance} \\
Geometry and semantics unresolvable in opposite places under a paired Student-t intervals & \textbf{Ours, claim-grade} & Phase~A, Table~\ref{tab:phasea} \\
RankMe and participation ratio each blind to the mode the other detects & \textbf{Ours, observed} & Phase~A, Table~\ref{tab:rank}; mechanism is centered vs.\ raw spectrum \\
Diagnostic reliability transferring to scenario retrieval & \textbf{Ours, in progress} & Phase~B; CIFAR-10 calibration complete \\
Projector as the variable governing where a distributional regularizer may act & \textbf{Ours, tested} & Survives its registered geometric test, gains no semantic support \\
\bottomrule
\end{tabular}
\end{table}

Two caveats are load-bearing and are repeated wherever the corresponding results appear.

First, the scale-dependence result this report attempts to replicate \replication{} was
encountered through a secondhand summary rather than a full reading of its source. It
should be treated as unverified.

Second, our SIGReg implementation was written from a summary-level understanding of
\citet{lejepa2025} and, when finally checked against the reference implementation, was
found to be \emph{wrong} in a way that silently disabled the regularizer: we centered the
embeddings before the normality test, which the reference does not, making a collapsed
representation a stationary point of the objective meant to prevent collapse.
Section~\ref{sec:instrument} reports this. It has been corrected and now matches the
reference bit-for-bit on isotropic, anisotropic, collapsed and mean-shifted inputs. We note
it prominently because it is the clearest possible vindication of labelling unvalidated
work as unvalidated: the implementation passed every behavioural test we had thought to
write, and was still wrong. \interp{}

A third load-bearing caveat was added after an external audit (2026-08-14): every SIGReg
result collected before Phase 2 comes from a configuration with \emph{no projector}, i.e.
the regularizer acted directly on the probed representation. The literature reports the
projector is worth on the order of twenty accuracy points in comparable settings, so those
results characterise the no-projector configuration only — not SIGReg, and not
\citet{lejepa2025}. The audit also found the original condition grid could not isolate
stop-gradient (no \texttt{none\_stopgrad} cell) and that no experiment had been repeated
across seeds; both are corrected in the pre-registered Phase-2 design.
